% Source PDF: source/普通高考/2012/2012湖北文.pdf, page 1, question 16.
% pdfimages -list reports no embedded bitmap; the flowchart is vector/text artwork.
% Grid evidence: tmp/review_2012_hubei_liberal/grid/page1_grid100.png (100px coarse grid).
% No-grid source crop: tmp/review_2012_hubei_liberal/crops/q16_source_tight.png.
% Nodes: 开始; a=1,s=0,n=1; s=s+a; a=a+2; n<3?; n=n+1; 输出s; 结束.
% Directed edges: 开始→初始化→s更新→a更新→判定;
% 判定(是)→n更新→回线→共享入口→s更新; 判定(否)→输出→结束.
% The return arrow points right into the shared entry; the main downward arrow
% into s=s+a is unique. The output branch is independent of the loop.

\begin{tikzpicture}[
  x=1cm,y=1cm,baseline,
  flowarrow/.style={-{Latex[length=2.1mm,width=1.5mm]},line width=.45pt},
  nodebox/.style={draw,line width=.45pt,inner sep=2pt,minimum height=.68cm,
    anchor=center,align=center,execute at begin node={\setlength{\parindent}{0pt}\noindent}},
  branchlabel/.style={anchor=center,align=center,inner sep=0pt,
    execute at begin node={\setlength{\parindent}{0pt}\noindent}},
  term/.style={nodebox,rounded corners=2.5pt,minimum width=1.30cm,text width=1.30cm},
  io/.style={nodebox,shape=trapezium,trapezium left angle=70,trapezium right angle=110,
    minimum width=2.30cm,text width=1.90cm,minimum height=.72cm},
  decision/.style={nodebox,diamond,aspect=2.7,minimum width=3.55cm,minimum height=1.02cm,inner sep=1pt}
]
  \node[term] (start) at (0,10.20) {开始};
  \node[nodebox,text width=4.15cm,minimum width=4.15cm] (init) at (0,8.85)
    {$a=1,s=0,n=1$};
  \node[nodebox,text width=2.55cm,minimum width=2.55cm] (sum) at (0,7.45)
    {$s=s+a$};
  \node[nodebox,text width=2.55cm,minimum width=2.55cm] (increment) at (0,6.10)
    {$a=a+2$};
  \node[decision] (test) at (0,4.65) {$n<3?$};
  \node[nodebox,text width=2.50cm,minimum width=2.50cm] (next) at (0,3.00)
    {$n=n+1$};
  \node[io] (output) at (4.10,3.95) {输出 $s$};
  \node[term] (finish) at (4.10,2.55) {结束};

  \coordinate (loop-left) at (-3.05,8.12);
  \coordinate (loop-right) at (4.10,8.12);
  \coordinate (loop-drop) at (0,2.35);
  \coordinate (loop-bottom) at (-3.05,2.35);

  % Initialization stem and the single downward arrow into s=s+a.
  \draw[flowarrow] (start.south) -- (init.north);
  % The return arrow lands on the interior of the main stem; avoid a
  % parsed 3-way endpoint while preserving the original loop topology.
  \draw[flowarrow] (init.south) -- (sum.north);
  \draw[flowarrow] (sum.south) -- (increment.north);
  \draw[flowarrow] (increment.south) -- (test.north);

  % Yes branch continues to n=n+1 and returns around the left side.
  \draw[flowarrow] (test.south) -- node[right=2pt] {是} (next.north);
  \coordinate (loop-entry) at (0,8.20);
  \draw (next.south) -- (loop-drop) -- (loop-bottom) -- (loop-left);
  \draw[flowarrow] (loop-left) -- (loop-entry);

  % No branch exits to output and termination, independent of the loop.
  \draw[flowarrow] (test.east) -- (4.10,4.65) -- (output.north);
  \draw[flowarrow] (output.south) -- (finish.north);
  \node[branchlabel] at (2.10,5.08) {否};

  \path[use as bounding box] (-3.40,2.00) rectangle (5.35,10.60);
\end{tikzpicture}
